% pannon-header.tex – a Pannon tankönyvsorozat PDF-arculata (XeLaTeX)

% A keretes megjegyzések (callout) ikon nélkül jelennek meg – a sorozat visszafogott,
% piktogrammentes arculata szerint. A fontawesome5 csomagot betöltöttnek jelöljük, az
% ikonparancsok üresek (így a fordítás a csomag telepítése nélkül is működik).
\makeatletter
\@namedef{ver@fontawesome5.sty}{2026/01/01}
\makeatother
\providecommand{\faInfo}{}\providecommand{\faExclamation}{}\providecommand{\faLightbulb}{}
\providecommand{\faExclamationTriangle}{}\providecommand{\faFire}{}
\newcommand{\pecalloutgap}{}
\usepackage{xcolor}
\definecolor{peblue}{HTML}{285780}
\definecolor{pebluedark}{HTML}{1C3F5E}
\definecolor{pelight}{HTML}{E8EEF4}
\definecolor{peorange}{HTML}{C8711E}
\definecolor{pegreen}{HTML}{1F9077}
\definecolor{peink}{HTML}{1F2933}
\definecolor{pemuted}{HTML}{5B6775}

\usepackage{fontspec}
\newfontfamily\pefontheading{Raleway}[Path=assets/fonts/, Extension=.otf, Ligatures=TeX,
  UprightFont=*-Bold, BoldFont=*-Black]
\newfontfamily\pefontblack{Raleway}[Path=assets/fonts/, Extension=.otf, Ligatures=TeX,
  UprightFont=*-Black, BoldFont=*-Black]
\newfontfamily\pefontsans{IBMPlexSans}[Path=assets/fonts/, Extension=.otf, Ligatures=TeX,
  UprightFont=*-Regular, BoldFont=*-SemiBold, ItalicFont=*-Italic]

% KOMA-Script címsorok
\setkomafont{disposition}{\pefontheading\color{pebluedark}}
\setkomafont{chapter}{\pefontblack\Huge\color{pebluedark}}
\setkomafont{chapterprefix}{\pefontheading\large\color{peblue}}
\setkomafont{part}{\pefontblack\color{peblue}}
\setkomafont{partnumber}{\pefontheading\color{peblue}}
\setkomafont{caption}{\pefontsans\small}
\setkomafont{captionlabel}{\pefontsans\small\bfseries\color{peblue}}
\setkomafont{pageheadfoot}{\pefontsans\footnotesize\color{pemuted}}

\usepackage{tikz}
% ferde sáv a fejezetcím alatt – a sorozat azonosító eleme
\makeatletter
\renewcommand*{\chapterlinesformat}[3]{%
  \@hangfrom{#2}{#3}\par\nobreak\vspace{.6ex}%
  \tikz\fill[peblue] (0,0) -- (3.4,0) -- (3.2,-.22) -- (-.2,-.22) -- cycle;}
\makeatother

% ábra- és táblázatfeliratok magyar sorrendben: "1.1. ábra."
\usepackage{caption}
\DeclareCaptionLabelFormat{hu}{#2.~#1}
\captionsetup{labelformat=hu, labelsep=period, font=small, justification=raggedright,
  singlelinecheck=false}

% didaktikai dobozok
\usepackage[most]{tcolorbox}
\newtcolorbox{peroadmap}{enhanced, breakable, colback=pelight, colframe=pelight,
  boxrule=0pt, arc=2pt, left=8pt, right=8pt, top=6pt, bottom=6pt,
  fontupper=\pefontsans\small,
  title={Miről szól a fejezet?}, coltitle=peblue, fonttitle=\pefontsans\bfseries\small,
  attach title to upper={\par\smallskip}}
\newtcolorbox{peesetuzleti}{enhanced, breakable, colback=peorange!7, colframe=peorange,
  boxrule=0pt, leftrule=3pt, arc=0pt, left=8pt,
  title={FUTÓ ESET · PANNONSHOP}, coltitle=peorange, fonttitle=\pefontsans\bfseries\footnotesize,
  attach title to upper={\par\smallskip}}
\newtcolorbox{peesetmernoki}{enhanced, breakable, colback=pegreen!7, colframe=pegreen,
  boxrule=0pt, leftrule=3pt, arc=0pt, left=8pt,
  title={FUTÓ ESET · PANNON ÜZEM}, coltitle=pegreen, fonttitle=\pefontsans\bfseries\footnotesize,
  attach title to upper={\par\smallskip}}

% unicode-math mellett a \boldsymbol a félkövér dőlt matematikai betűket használja (a bm csomag ütközne)
\AtBeginDocument{\renewcommand{\boldsymbol}[1]{\symbfit{#1}}\providecommand{\bm}[1]{\symbfit{#1}}}
\usepackage{mathtools}
\frenchspacing

% részcímek: a rész sorszáma a címben szerepel (pl. "I. Alapok"), ezért a KOMA-számozást elhagyjuk
\renewcommand*{\partformat}{}
\AtBeginDocument{%
  \renewcommand*{\figurename}{ábra}%
  \renewcommand*{\tablename}{táblázat}%
  \renewcommand*{\contentsname}{Tartalomjegyzék}%
  \renewcommand*{\listfigurename}{Ábrák jegyzéke}%
  \renewcommand*{\listtablename}{Táblázatok jegyzéke}%
  \renewcommand*{\appendixname}{Függelék}%
  \renewcommand*{\chaptername}{fejezet}}

% Az IBM Plex betűkészletekből hiányzó néhány jel matematikai megfelelője (szövegben és kódkimenetben is)
\usepackage{newunicodechar}
\newunicodechar{⇒}{\ensuremath{\Rightarrow}}
\newunicodechar{∈}{\ensuremath{\in}}
\newunicodechar{σ}{\ensuremath{\sigma}}
\newunicodechar{α}{\ensuremath{\alpha}}
\newunicodechar{λ}{\ensuremath{\lambda}}
\newunicodechar{μ}{\ensuremath{\mu}}
\newunicodechar{₂}{\ensuremath{_2}}
\newunicodechar{₄}{\ensuremath{_4}}
% a számított (pandas) táblázatok kisebb betűvel és szűkebb oszlopközzel, hogy elférjenek a B5-ös oldalon
\usepackage{etoolbox}
\AtBeginEnvironment{longtable}{\scriptsize\setlength{\tabcolsep}{3pt}}

% A hosszú kódsorok és kódkimenetek tördelése a szedéstükrön belül
\usepackage{fvextra}
\makeatletter
\AtBeginDocument{%
  \@ifundefined{Highlighting}{}{%
    \RecustomVerbatimEnvironment{Highlighting}{Verbatim}{breaklines,breakanywhere,commandchars=\\\{\},fontsize=\small}}%
  \RecustomVerbatimEnvironment{verbatim}{Verbatim}{breaklines,breakanywhere,fontsize=\small}}
\makeatother
